\documentclass[10pt,a4paper]{amsart}
\usepackage[margin=19mm]{geometry}
\usepackage{amsmath,amssymb,mathtools}
\usepackage[scr=rsfs,cal=euler]{mathalfa}
\usepackage{xfrac}
\usepackage[unicode,hidelinks,bookmarks=false]{hyperref}
\newcommand{\ie}{i.e.\ }
\newcommand{\cf}{cf.\ }
\newcommand{\resp}{respectively\ }
\newcommand{\Subsection}{\subsection}
\providecommand{\nobreakdash}{\nobreak\hbox{-}}
\setlength{\emergencystretch}{2em}
\multlinegap0pt
\usepackage{bm}
\usepackage{dsfont}
\usepackage{algorithm}\usepackage{algpseudocode}
\usepackage{natbib}
\newtheorem{theorem}{Theorem}
\newtheorem{lemma}{Lemma}
\usepackage{cleveref}
\crefname{theorem}{Theorem}{Theorems}
\crefname{lemma}{Lemma}{Lemmas}
\newcommand{\E}{{\mathbb{E}}}
\newcommand{\var}{\operatorname{Var}}
\title{Bandit Allocational Instability}
\author{Yilun Chen, Jiaqi Lu}
\date{Source: arXiv:2602.07472v1 (CC BY 4.0)}
\begin{document}
\maketitle

\noindent\emph{Source and licence.} Bandit Allocational Instability. Version arXiv:2602.07472v1 (CC BY 4.0), distributed by arXiv under the Creative Commons Attribution 4.0 International licence, http://creativecommons.org/licenses/by/4.0/, as recorded in the arXiv licence metadata for this exact version and shown on the abstract page https://arxiv.org/abs/2602.07472v1. Full licence notice: \texttt{licenses/arxiv-2602.07472-CC-BY-4.0.txt}.\par\smallskip

\noindent\textit{Editorial scope.} Counted unit: the article's theorem thm:ucb-f (source0.tex lines 328--335). The article's complete original proof of it is delivered with it (source0.tex lines 915--972, one \texttt{proof} environment of 58 lines, which the article places away from the statement). No mathematics is written, repaired or re-derived: the statement and the proof are the article's own lines, in the article's own notation, and no retained line is substituted or re-flowed.  Retained uncounted as the source-local context the proof needs: lines 725--727; lines 734--747.  Nothing was omitted from any retained span, so the package carries no omission record.  No retained line contains a citation, so the wrapper carries no bibliography and no bibliography entry.  No retained line contains a figure environment, an image-inclusion command or drawing code, so no image is delivered and the package declares no asset dependency.  Editorial formatting only, in addition: an AMS \texttt{amsart} preamble that restates the article's own declarations and macros used by the retained lines, namely the environments \texttt{theorem}, \texttt{lemma} on separate counters, together with the standard packages those lines need.  The retained environments print in this document as Theorem 1, Lemma 1 in order of appearance, whereas the article prints the same items with the numbering of its own full document.  The complete native source of the article as distributed in the arXiv e-print for this exact version is bundled unchanged as \texttt{sources/194182/source0.tex}.
\begin{lemma}\label{lem:n_i_asymptotics}
 If the exploration function $f(\cdot)$ satisfies both $f(t)$ and $\frac{\sqrt{t}}{f(t)}$ increasing in $t$, then $\inf_{\nu\in\mathfrak{P}_{\text{sg}}}n_i$ and $\inf_{\nu\in\mathfrak{P}_{\text{sg}}}\frac{n_2}{f(T)}$ both diverges as $T \to \infty$ for any $i=2,...,K$. Moreover, for such $f(\cdot)$, $\sup_{\nu\in\mathfrak{P}_{\text{sg}}}(n_i\Delta_i)=O\left(T^{\frac{1}{2}}f(T)\right)$ for any $i=2,...,K$. 
\end{lemma}

\begin{lemma}\label{lem:N_i_tail_bound}
Consider \texttt{UCB-f} with an exploration function $f(\cdot)$ such that both
$f(t)/\log t$ and $\sqrt{t}/f(t)$ are non-decreasing in $t$. There exists a constant
$\underline{T}\equiv \underline{T}(M,\sigma,K,\alpha)$ such that, for any
$\nu\in\mathfrak{P}_{\mathrm{sg}}$, any $T\ge \underline{T}$, and any $i\in\{1,\dots,K\}$,
\begin{align*}
P_{\nu}\!\left(N_i-n_i>q\right)
&\le \frac{K-1}{4q^4},
\qquad \forall\, q\ge \frac{\log T}{f(T)}\left(n_i\wedge n_2\right),\\
P_{\nu}\!\left(N_i-n_i<-q\right)
&\le \frac{(K-1)^6}{4q^4},
\qquad \forall\, q\ge \frac{(K-1)\log T}{f(T)}\,n_2.
\end{align*}
\end{lemma}

\begin{theorem}\label{thm:ucb-f}
For any exploration function $f(\cdot)$ such that both $f(t)/\log t$ and $\sqrt{t}/f(t)$ are non-decreasing in $t$, \texttt{UCB-f} achieves, over the instance class $\mathfrak{P}_{\mathrm{sg}}$,
\[
\mathcal{R}_T = O\!\left(T^{\frac12}f(T)\right),
\qquad
\mathcal{S}_T = O\!\left(\frac{T \log T}{f(T)}\right).
\]
\end{theorem}

\begin{proof}{Proof of \Cref{thm:ucb-f}}
Fix $\alpha\in[0,\tfrac12)$. Recall $n_i = n_i^T$ is the solution of the fluid system at horizon $T$, and we use $N_i = N_{i,T}$ interchangeably.
\paragraph{Allocation variability.}
By \Cref{lem:N_i_tail_bound}, for any $T\ge \underline T$, any $\nu \in \mathfrak{P}_{\text{sg}}$ and any $i\in[K]$,
\[
P\big(|N_i-n_i|>q\big)\ \le\ \frac{(K-1)^6}{2q^4},
\qquad \forall q\ge \frac{(K-1)\log T}{f(T)}n_2.
\]
Fix $q_0\triangleq \frac{(K-1)\log T}{f(T)}n_2$. Using
$\E[Z^2]=\int_0^\infty 2u\,P(|Z|>u)\,du$, we obtain
\begin{align*}
\E\big[(N_i-n_i)^2\big]
&= \int_0^{q_0} 2u\,P\big(|N_i-n_i|>u\big)\,du
   +\int_{q_0}^{\infty} 2u\,P\big(|N_i-n_i|>u\big)\,du \\
&\le q_0^2 + \int_{q_0}^{\infty} 2u\cdot \frac{(K-1)^6}{2u^4}\,du
 \;=\; q_0^2 + \frac{(K-1)^6}{2}\int_{q_0}^{\infty} u^{-3}\,du \\
&= \frac{(K-1)^2 (\log T)^2 n_2^2}{f(T)^2}+\frac{(K-1)^4 f(T)^2}{4n_2^2(\log T)^2}.
\end{align*}
By \Cref{lem:n_i_asymptotics}, $\inf_{\nu\in\mathfrak{P}_{\text{sg}}} n_2/f(T)\to\infty$ as $T\to\infty$.
Hence there exists $T_0=T_0(K,\alpha)$ such that for all $T\ge T_0$ and all $\nu\in\mathfrak{P}_{\text{sg}}$,
\[
\frac{(K-1)^4 f(T)^2}{4n_2^2(\log T)^2}\ \le\ \frac{(K-1)^2 (\log T)^2 n_2^2}{f(T)^2}.
\]
Therefore, for all $T\ge \underline{\underline T}\triangleq \underline T\vee T_0$, $\E\big[(N_i-n_i)^2\big]\ \le\ 2\frac{(K-1)^2 (\log T)^2 n_2^2}{f(T)^2},$ and hence 
\[
\sqrt{\var(N_i)}\ \le\ \sqrt{\E[(N_i-n_i)^2]}\ \le\ \sqrt{2}(K-1)\frac{\log T}{f(T)}n_2.
\]
Since $n_2=O(T)$, we obtain 
\[
\mathcal S_T
=\sup_{\nu\in\mathfrak{P}_{\text{sg}}}\max_{i\in[K]}\sqrt{\var_\nu(N_{i,T})}
= O\left(\frac{T\log T}{f(T)}\right).
\]

\paragraph{Regret.}
We bound $\E N_i$ via the right tail. By \Cref{lem:N_i_tail_bound}, for $T\ge \underline T$ and
any $i\in[2, K]$,
\[
P(N_i-n_i>q)\ \le\ \frac{K-1}{4q^4},
\qquad \forall q\ge \frac{\log T}{f(T)}n_i.
\]
Taking $q=n_i$ and using $\E[Y]=\int_0^\infty P(Y>u)\,du$ for $Y\ge 0$, we get
\begin{align*}
\E N_i
&\le n_i + \E[(N_i-n_i)_+]
 \le n_i + \int_0^{n_i} 1\,du + \int_{n_i}^{\infty} P(N_i-n_i>u)\,du \\
&\le 2n_i + \int_{n_i}^{\infty} \frac{K-1}{4u^4}\,du  \;=\;2n_i + \frac{K-1}{12\,n_i^3}.
\end{align*}
Since by Lemma~\ref{lem:n_i_asymptotics}, $\inf_{\nu\in\mathfrak{P}_{\text{sg}}}n_i\to\infty$, there exists $T_0'=T_0'( M, K)$ such that $\frac{K-1}{12n_i^3}\le 2n_i$ for all $i$ and all $\nu$ whenever $T\ge T_0'$. Thus for $T\ge \underline T\vee T_0'$, $\E_{\nu} N_i \le 4n_i,$ for all $i\in[2, K]$ and $\nu \in \mathfrak{P}_{\text{sg}}$ Therefore,
\[
\sup_{\nu\in\mathfrak{P}_{\text{sg}}} R_T(\nu)
=\sup_{\nu\in\mathfrak{P}_{\text{sg}}}\sum_{i=2}^K \Delta_i\,\E_{\nu} N_i
\le 4\sup_{\nu\in\mathfrak{P}_{\text{sg}}}\sum_{i=2}^K \Delta_i\,n_i
\le 4\sum_{i=2}^K \sup_{\nu\in\mathfrak{P}_{\text{sg}}}(n_i\Delta_i)
=O\!\left(T^{\frac12}f(T)\right),
\]
where the last step again uses \Cref{lem:n_i_asymptotics}. This completes the proof. \qed
\end{proof}

\end{document}
